\section{The endomorphism \texorpdfstring{$\Omega$}{Omega}, the flow and the function \texorpdfstring{$\mathcal S_{\mathbb C}$}{S}}\label{sec:setting}

This section fixes notation and recalls what is used from \cite{hkkp26} and
\cite{wue26}. Nothing in it is new. Throughout, $T^\dagger$ is the adjoint of
a linear map for the Hermitian metrics in force (the conjugate transpose in
orthonormal frames), $T^*$ is the conjugate transpose of a matrix in fixed
reference bases, and $D^+$ is the Moore--Penrose inverse. Inner products
$\inner\cdot\cdot$ and Hermitian forms are conjugate-linear in the first
argument, and a scalar in a matrix expression denotes the corresponding
multiple of the identity. For a square matrix
$T$ we write
$\Rea T=(T+T^\dagger)/2$ and $\Ima T=(T-T^\dagger)/2\ii$; $\norm T_1$ is the
trace norm, $\norm T_{\HS}$ the Hilbert--Schmidt norm and $\norm T$ the
operator norm; for a Hermitian matrix $M$, $M_\pm\geq0$ are its positive and
negative parts, $M=M_+-M_-$, $M_+M_-=0$. The letter $\varphi$ always denotes
the argument of the total central charge; $\phi$ with indices denotes phases
of blocks and their eigenvalues. Likewise $\vartheta$, $\varkappa$,
$\varrho$, $\varsigma$, $\varpi$ denote objects different from those denoted
by $\theta$, $\kappa$, $\rho$, $\sigma$, $\pi$.

\subsection{Complexes and the endomorphism \texorpdfstring{$\Omega$}{Omega}}
Let $V,U,W$ be finite-dimensional complex vector spaces, $D:V\to U$ a linear
map and $A_1,\dots,A_k:V\to W$ linear maps. This is a two-term complex of
representations of the generalized Kronecker quiver, with two vertices $0,1$
and $k$ arrows from $0$ to $1$:
\begin{equation}\label{eq:complex}
 \mathsf E:\qquad(V\xrightarrow{\,(A_j)\,}W)\ \xrightarrow{\,(D,0)\,}\ (U\to0),
\end{equation}
with $V$ and $W$ in degree $0$ and $U$ in degree $1$. We assume throughout
that
\[
 D\text{ is surjective},\qquad n_0:=\dim\ker D\geq1,\quad n_1:=\dim W\geq1,
 \quad r:=\dim U\geq1 .
\]
The \emph{cohomology representation} is $\ker D\to W$ with the arrows
$A_j|_{\ker D}$, of dimension vector $(n_0,n_1)$. It is a subcomplex of
$\mathsf E$, and the quotient $(V/\ker D\to0)\to(U\to0)$ is an acyclic
complex, which we call the \emph{acyclic part}; $r$ is its rank. Fix
\emph{central charges} $z_0,z_1$ in the open upper half-plane, attached to
the two vertices, and write
\[
 z_i=x_i+\ii y_i=|z_i|e^{\ii\alpha_i},\qquad y_i>0,\quad\alpha_i\in(0,\pi).
\]

Let $h=(h_V,h_U,h_W)$ be Hermitian metrics. Adjoints refer to the current
metrics. Put
\[
 B=D^\dagger D,\qquad L=DD^\dagger,\qquad C=\sum_jA_j^\dagger A_j,
\]
and, for Hermitian matrices $P\geq0$ and $M$,
\begin{equation}\label{eq:J}
 \begin{gathered}
 R(P)=\int_0^1e^{-xP}dx,\qquad
 K_P(M)=\int_0^1e^{-(1-x)P}Me^{-xP}dx,\\
 J_P(M)=\int_{\substack{x,y\geq0\\x+y\leq1}}e^{-xP}Me^{-yP}\,dx\,dy .
 \end{gathered}
\end{equation}
In \cite[\S5.4]{hkkp26} a complex of quiver representations with metrics
carries an endomorphism $\Omega$ of its underlying graded vector space,
defined through a superconnection in a graded algebra attached to the
quiver: $\Omega=\varepsilon\,[\exp(-\sqrt{-1}F)]_f$, where
$F=(A-A^*)^2+\sqrt{-1}\sum_if_i\otimes z_i$ is formed from the element
$A=\sum_ie_i\otimes d_i+\sum_aa\otimes\rho(a)$ of the algebra of
\cite[Definition~5.13]{hkkp26} tensored with the endomorphisms, with $d_i$
the differentials and $\rho(a)$ the maps of the arrows, $\varepsilon$ is the
grading, and $[\,\cdot\,]_f$ is the component along the basis elements
$f_i$. Apart from $\Omega$ and $z_i$, the symbols of this sentence are those
of \cite{hkkp26} and have this meaning only here; $A^*$ is formed from the
involution of that algebra and the adjoints for the metrics
\cite[(5.14)]{hkkp26}. The
differential of the function $\cS_{\CC}$ below is
$\xi\mapsto\tr(\Omega\xi)$; the trace of $\Omega$ is the central charge
\eqref{eq:central}, and for a representation $\Omega$ is the central charge
of each vertex shifted by the moment map. For \eqref{eq:complex} it is
$\Omega=\Omega_V\oplus\Omega_U\oplus\Omega_W$ with
\cite[Proposition~2.1]{wue26}
\begin{equation}\label{eq:coeff}
 \Omega_V=z_0e^{-B}+K_B(C),\qquad
 \Omega_U=J_L(DCD^\dagger)-z_0e^{-L},\qquad
 \Omega_W=z_1\id-\sum_jA_jR(B)A_j^\dagger .
\end{equation}
We abbreviate
\begin{equation}\label{eq:KJG}
 K=K_B(C),\qquad J=J_L(DCD^\dagger),\qquad G=\sum_jA_jR(B)A_j^\dagger\geq0 ,
\end{equation}
three Hermitian matrices that do not depend on the charges. Hence
\begin{equation}\label{eq:imaginary}
 \Ima\Omega_V=y_0e^{-B}>0,\qquad
 \Ima(-\Omega_U)=y_0e^{-L}>0,\qquad
 \Ima\Omega_W=y_1\id>0,
\end{equation}
and \cite[Lemma~5.17]{hkkp26}
\begin{equation}\label{eq:central}
 \tr\Omega_V+\tr\Omega_U+\tr\Omega_W=n_0z_0+n_1z_1=:Z .
\end{equation}
The matrix $\Omega_W$ is normal. The matrix $\Omega_V$ is in general not
normal \cite[Example~2.3]{wue26}, and neither is $\Omega_U$, whose Hermitian
and anti-Hermitian parts $J-x_0e^{-L}$ and $-y_0e^{-L}$ commute only if $J$
commutes with $e^{-L}$.

\subsection{The phase and the flow}
For a matrix $T$ with $\Ima T>0$ there is a unique Hermitian matrix $\Psi$
with spectrum in $(0,\pi)$ such that $Te^{-\ii\Psi}>0$, the \emph{polar phase}
of $T$; it is the argument of the unitary factor in the polar decomposition
$T=(TT^\dagger)^{1/2}e^{\ii\Psi}$ \cite[Section~7.3]{hj13}: if
$e^{\ii\Psi}v=\zeta v$, then
$\inner v{(\Ima T)v}=(\Ima\zeta)\inner v{|T|v}$, so the eigenvalues of the
unitary factor lie in the open upper half-circle (see also
\cite[Lemma~2.2]{wue26}). Here
$|T|:=(TT^\dagger)^{1/2}$ is the left modulus; the unitary factor is the same
in the left and right polar decompositions, and so is the trace norm (in
\cite{bha97} the modulus is the right one, $(T^\dagger T)^{1/2}$). The phases
$\Phi_V,\Phi_W$ are the polar phases of $\Omega_V,\Omega_W$, and
$\Phi_U+\pi$ is the polar phase of $-\Omega_U$; so
\begin{equation}\label{eq:phase}
 |\Omega_\bullet|:=\Omega_\bullet e^{-\ii\Phi_\bullet}
 =(\Omega_\bullet\Omega_\bullet^\dagger)^{1/2}>0,\qquad\bullet\in\{V,U,W\},
\end{equation}
with $\spec\Phi_V,\spec\Phi_W\subset(0,\pi)$ and
$\spec\Phi_U\subset(-\pi,0)$. This is the phase defined by positivity in
\cite[\S4.2]{hkkp26}, with the branches of \cite[(5.15)]{hkkp26}. In
\cite[\S5.4]{hkkp26} the phase of $\Omega$ for a complex is written with the
functional calculus. For normal $\Omega$, in particular for
representations, the two readings coincide; for complexes $\Omega$ is not
normal in general \cite[Example~2.3]{wue26}. The flow studied here is the
one defined with the polar phase:
\begin{equation}\label{eq:flow}
 h_\bullet^{-1}\dot h_\bullet=-\Phi_\bullet,\qquad\bullet\in\{V,U,W\}.
\end{equation}
We call it the \emph{phase flow}. It is the minimizing flow of
\cite[\S4.1]{hkkp26} with the phase read as above; we use a neutral name
because on complexes the mass can increase along it
\cite[Theorem~1.1]{wue26}. The \emph{mass} is $m(h)=\sum_\bullet\tr|\Omega_\bullet|\geq|Z|$. The
vector field is smooth and its solutions exist for all $t\geq0$
\cite[Lemma~2.4]{wue26}. In orthonormal frames moving with the metrics the
maps evolve by
\begin{equation}\label{eq:mapflow}
 \dot D=\tfrac12(D\Phi_V-\Phi_UD),\qquad
 \dot A_j=\tfrac12(A_j\Phi_V-\Phi_WA_j).
\end{equation}
A \emph{state} is a tuple $(D,A_j)$ of matrices in orthonormal frames, up to
unitary changes of the frames; equivalently, a triple of metrics on the
fixed complex up to isometry, or a point of the orbit of $(D,A_j)$ under
$\mathrm{GL}(V)\times\mathrm{GL}(U)\times\mathrm{GL}(W)$ modulo the unitary
groups. The phases, mass and arrow energy are invariant under these changes of frames. The primitive below also uses fixed reference volume elements.

\begin{lemma}\label{lem:eigenbasis}
Let $T$ be a matrix with $\Ima T>0$ and polar phase $\Psi$, and let $v$ be a
unit vector with $\Psi v=\psi v$. Then
$\inner v{Tv}=p\,e^{\ii\psi}$ with $p=\inner v{|T|v}>0$. Consequently, in an
orthonormal eigenbasis $(v_k)$ of $\Psi$ with eigenvalues $\psi_k$ and
$p_k=\inner{v_k}{|T|v_k}$, $\tr(|T|f(\Psi))=\sum_kp_kf(\psi_k)$ for every
real function $f$.
\end{lemma}

\begin{proof}
$T=|T|e^{\ii\Psi}$, so $\inner v{Tv}=e^{\ii\psi}\inner v{|T|v}$. The second
statement is the trace computed in the basis $(v_k)$.
\end{proof}

The lemma applies to $\Omega_V$, $\Omega_W$ and $-\Omega_U$; for the block
$U$ it gives $\inner v{\Omega_Uv}=p\,e^{\ii\phi}$ for a unit eigenvector $v$
of $\Phi_U$ with eigenvalue $\phi\in(-\pi,0)$.

\subsection{The primitive}
Let $\Ein(b)=\int_0^1(1-e^{-xb})\,dx/x$. With respect to fixed reference
bases, in which $\log\det h_\bullet$ denotes the logarithm of the Gram
determinant, put
\begin{equation}\label{eq:primitive}
 \cS_{\CC}(h)=z_0\bigl(\log\det h_V-\log\det h_U\bigr)+z_1\log\det h_W
 +z_0\tr\Ein(B)-\tr\bigl(CR(B)\bigr).
\end{equation}
Then \cite[Proposition~7.1]{wue26}, for every variation
$\xi_\bullet=h_\bullet^{-1}\delta h_\bullet$ of the metrics,
\begin{equation}\label{eq:primitivevariation}
 d\cS_{\CC}(\xi)=\sum_\bullet\tr(\Omega_\bullet\xi_\bullet).
\end{equation}
Both sides are affine in $(z_0,z_1)$, so the terms without charges agree:
\begin{equation}\label{eq:energyvariation}
 d\bigl(\tr(CR(B))\bigr)(\xi)=-\tr(K\xi_V)-\tr(J\xi_U)+\tr(G\xi_W).
\end{equation}

\subsection{Charges and the stability parameter}
Put
\[
 \varphi=\arg Z\in(0,\pi),\quad\sigma=\sin\varphi,\quad
 c_0=|z_0|\sin(\alpha_0-\varphi),\quad c_1=|z_1|\sin(\varphi-\alpha_1).
\]
Since $\Ima(e^{-\ii\varphi}Z)=0$ we have $n_0c_0=n_1c_1$. The charges are
\emph{ordered} if $\alpha_0>\alpha_1$ and \emph{reversed} if
$\alpha_0<\alpha_1$; they are ordered exactly when $c_0>0$, or when
$\alpha_1<\varphi<\alpha_0$. For ordered charges put
\begin{equation}\label{eq:theta}
 \theta_0=c_0/\sigma,\qquad\theta_1=c_1/\sigma,\qquad
 \vartheta=n_0\theta_0=n_1\theta_1>0 ;
\end{equation}
then $z_0+\theta_0$ and $z_1-\theta_1$ have argument $\varphi$. We call the
pair $(\theta_0,\theta_1)$ the \emph{stability parameter}. A
representation whose dimension vector $(a_0,a_1)$ satisfies
$a_0\theta_0=a_1\theta_1$, for instance $(n_0,n_1)$, has phase $\varphi$. It
is \emph{semistable} if $a\theta_0-b\theta_1\leq0$ for the dimension vector
$(a,b)$ of each of its subrepresentations, that is, if no subrepresentation
has phase $\arg(az_0+bz_1)$ larger than $\varphi$; it is \emph{stable} if the
inequality is strict for proper nonzero subrepresentations, and
\emph{polystable} if it is a direct sum of stable representations of phase
$\varphi$. In the convention of King \cite[Definition~1.1]{kin94} this is
stability for the character $(-\theta_0,\theta_1)$, a positive multiple of
the integral character $(-n_1,n_0)$; the Kempf--Ness function \eqref{eq:KN}
carries the weights $(\theta_0,-\theta_1)$ on
$(\log\det h_0,\log\det h_1)$. For reversed charges no
representation with $n_0,n_1\geq1$ is semistable, because the
subrepresentation supported at the vertex $1$ has phase
$\alpha_1>\varphi$.

\begin{definition}\label{def:admissible}
Let $(n_0,n_1,r)$ be given. The charges $(z_0,z_1)$ are
\emph{admissible} if they are ordered and either $\varphi\geq\frac\pi3$ or
$r=1$.
\end{definition}

This is the range of charges covered by Theorem~\ref{thm:monotone}; the
number $\frac\pi3$ comes from the constant of Theorem~\ref{thm:kernel}.

The \emph{symmetric charges} are $z_0=-1+\ii$, $z_1=1+\ii$. If moreover
$n_0=n_1$, then $\varphi=\frac\pi2$, $\sigma=1$, $\theta_0=\theta_1=1$, and they
are admissible for every $r$.

\subsection{The splitting adapted to \texorpdfstring{$\ker D$}{ker D}}
Choose orthonormal frames adapted to $V=\ker D\oplus(\ker D)^\perp$ and write
\[
 D=[0,F],\qquad A_j=[H_j,T_j],\qquad
 \Phi_V=\begin{pmatrix}\Phi_{HH}&\Phi_{HA}\\\Phi_{AH}&\Phi_{AA}\end{pmatrix}.
\]
Here the indices $H$ and $A$ refer to $\ker D$ and to its orthogonal
complement, which carries the acyclic part. So $F$ is invertible and the
$H_j$ are the matrices of the arrows of the cohomology representation for the
induced metrics. Let
$D^+=D^\dagger(DD^\dagger)^{-1}$ be the Moore--Penrose inverse of $D$ and
\[
 X_j=A_jD^+=T_jF^{-1}:U\to W .
\]
These maps vanish exactly when the arrows vanish on $(\ker D)^\perp$. Put
\begin{equation}\label{eq:qslambda}
 C_0=\sum_jH_j^\dagger H_j,\quad C_1=\sum_jH_jH_j^\dagger,\quad
 q=\tr C_0=\tr C_1,\quad s=\sum_j\norm{X_j}_{\HS}^2,\quad
 \lambda=\lambda_{\min}(L),
\end{equation}
let $\ell_1\leq\dots\leq\ell_r$ be the eigenvalues of $L=FF^\dagger$, so
$\lambda=\ell_1$, and
\begin{equation}\label{eq:Y}
 Y=\sum_jX_j(\id-e^{-L})X_j^\dagger\geq0,\qquad G=C_1+Y,\qquad
 \tr Y\geq(1-e^{-\lambda})\,s .
\end{equation}
With the choice of the moving frames made in \cite[\S6.1]{wue26} (the frames
of $V$ are rotated so that the splitting is preserved),
\begin{equation}\label{eq:adapted}
 \dot F=\tfrac12(F\Phi_{AA}-\Phi_UF),\qquad
 \dot H_j=\tfrac12(H_j\Phi_{HH}-\Phi_WH_j),
\end{equation}
and
\begin{equation}\label{eq:adaptedcoeff}
 \Omega_V=\begin{pmatrix}z_0\id+C_0&E_R\\E_R^\dagger&z_0e^{-F^\dagger F}+K_{F^\dagger F}(C_T)\end{pmatrix},
 \qquad\Omega_W=z_1\id-C_1-Y,
\end{equation}
with $C_T=\sum_jT_j^\dagger T_j$ and $E_R=\sum_jH_j^\dagger T_jR(F^\dagger F)$.
For the cohomology representation with its induced metrics the
endomorphism $\Omega$ is
\[
 \Omega_H=(z_0\id+C_0)\oplus(z_1\id-C_1);
\]
it is normal, with phases the Hermitian matrices
$\phi_0=\arg(z_0\id+C_0)$ and $\phi_1=\arg(z_1\id-C_1)$, and mass
$m_H=\norm{z_0\id+C_0}_1+\norm{z_1\id-C_1}_1$. We write
\[
 \delta_0=\Phi_{HH}-\phi_0,\qquad\delta_1=\Phi_W-\phi_1,\qquad
 \mathcal E_V=\Omega_V-\diag\bigl(z_0\id+C_0,\ z_0e^{-F^\dagger F}\bigr).
\]
We say that \emph{the acyclic part cancels} along a trajectory if
$\lambda\to\infty$ and
$\norm{\Omega_V-(z_0\id+C_0)\oplus0}_1+\norm{\Omega_U}_1+\norm{\Omega_W-(z_1\id-C_1)}_1\to0$.
The name refers to the central charge, to which the two terms of the acyclic
part contribute $z_0r$ and $-z_0r$; the statement is that their
contributions to $\Omega$ itself tend to zero.

For ordered charges we write
$\mu=(\mu_0,\mu_1)=(C_0-\theta_0\id,\ C_1-\theta_1\id)$ and call it the
\emph{moment map} of the cohomology representation. The pair
$(-\mu_0,\mu_1)$ is the gradient of the Kempf--Ness function
(Lemma~\ref{lem:knvariation}), that is, up to a constant factor the moment
map of $\mathrm U(n_0)\times\mathrm U(n_1)$ in the usual convention, shifted
by King's character; we change the sign of the first component so that both
components read $C_i-\theta_i$. A pair of metrics on
$\ker D$ and $W$ with $\mu=0$ is a \emph{zero of the moment map}; such pairs
exist exactly when the cohomology representation is polystable
\cite[Proposition~6.5]{kin94}, which extends the theorem of Kempf and Ness
\cite{kn79}, and they are the harmonic metrics of \cite{hkkp26} for a
representation, since $\mu=0$ if and only if both phases $\phi_0,\phi_1$ equal
$\varphi$.

\subsection{Estimates for \texorpdfstring{$A_jD^+$}{A\_jD+} and for the phases}
We use the following facts from \cite{wue26}, valid for all charges in the
upper half-plane and all dimensions. Put $\kappa_0=\pi/(2e)$ and
$d_0=\min(n_0,r)$.

\begin{proposition}[{\cite{wue26}}]\label{prop:imported}
Along the phase flow and at every metric:
\begin{enumerate}
\item[(i)] $\lambda$ is nondecreasing along the flow. By the formula for the
 one-sided derivative of the least eigenvalue of a differentiable family of
 Hermitian matrices, its right derivative is the minimum of
 $\lambda\bigl(\inner v{\Phi_{AA}v}-\inner u{\Phi_Uu}\bigr)$ over the unit
 vectors $v$ with $F^\dagger Fv=\lambda v$, where $u=Fv/\sqrt\lambda$
 \cite[Lemma~4.3 and \S6.1]{wue26}.
\item[(ii)] $\norm{E_R}_{\HS}\leq\sqrt{qs/\lambda}$,
 $\norm{K_{F^\dagger F}(C_T)}_1\leq r\kappa_0s$, $\norm{J}_1\leq rs$ and
 $\norm{Y}_1\leq s$; hence
 $\norm{\mathcal E_V}_1\leq2\sqrt{d_0qs/\lambda}+r\kappa_0s$
 \cite[Lemma~4.2, \S6.1 and \S6.3]{wue26}.
\item[(iii)] $\dot s\leq\bigl(-\alpha_1+\frac{4\pi q}{y_0\lambda}\bigr)s$
 \cite[\S6.2]{wue26}.
\item[(iv)] If $\sqrt{qs/\lambda}\leq y_0/8$, then
 $\norm{\delta_0}_{\HS}\leq\frac{4\pi}{y_0}\sqrt{qs/\lambda}$; and
 $\norm{\delta_1}_{\HS}\leq s/y_1$ always
 \cite[Lemma~3.1 and \S6.2]{wue26}.
\item[(v)] Let $T_\circ=|T_\circ|\,e^{\ii\Psi_\circ}$ be an invertible
 matrix in polar form, with $\Psi_\circ$ Hermitian, and $T=T_\circ+M$ a
 matrix with $\Ima T>0$ and polar phase $\Psi$. If $v$ is a unit vector with
 $|T_\circ|v=pv$ and $\Psi_\circ v=\psi v$, where $p>0$ and $0<\psi<\pi$,
 then $|\inner v{(\Psi-\psi)v}|\leq\pi\sqrt{\norm M_1/p}$
 \cite[Lemma~3.2]{wue26}.
\end{enumerate}
\end{proposition}

In \cite{wue26} these estimates were used when the initial value of
$\lambda$ is large: then the acyclic part cancels, for all charges
\cite[Theorem~1.3]{wue26}, and for the symmetric charges the moment map of
the cohomology representation tends to zero if its Kempf--Ness function is
bounded below \cite[Theorem~1.4]{wue26}. Whether the assumption on the
initial value of $\lambda$ can be removed was left open there.
